% ============================================================
% Inhalt: Stabilität als Grundzustand
% Der Permanentmagnet als makroskopischer Beweis
% für diskrete Spin-Stabilität im FFGFT-Rahmen
% März 2026
% ============================================================

% Dirac-Notation
\providecommand{\ket}[1]{\left|#1\right\rangle}
\providecommand{\bra}[1]{\left\langle#1\right|}
\providecommand{\braket}[2]{\left\langle#1\middle|#2\right\rangle}

% ============================================================
\section{Die unbequeme Frage}
\label{sec:unbequeme-frage}
% ============================================================

In jedem Lehrbuch der Quantenmechanik steht:
Ein Elektron befindet sich in echter Superposition
$\alpha\ket{\uparrow} + \beta\ket{\downarrow}$ ---
gleichzeitig in beiden Spinzuständen,
bis eine Messung den Zustand kollabiert.

Auf jedem Kühlschrank klebt ein Permanentmagnet.

Zwischen diesen beiden Tatsachen liegt ein Widerspruch,
den die Standardtheorie nie explizit auflöst.

\begin{significance}[Die zentrale Frage dieses Dokuments]
Wenn ein einzelnes Elektron in echter Superposition
beider Spinzustände existiert ---
warum zeigen dann Milliarden von Elektronen
im Permanentmagneten \emph{niemals} kollektive Superposition,
sondern immer eindeutige, stabile Ausrichtung?

Und umgekehrt:
Was sagt der Permanentmagnet --- ein Alltagsgegenstand,
kein Laborexperiment --- über die Natur des Spins?
\end{significance}

% ============================================================
\section{Was der Permanentmagnet wirklich zeigt}
\label{sec:permanentmagnet}
% ============================================================

Ein Permanentmagnet ist kein mysteriöses Quantenphänomen.
Er ist die \textbf{makroskopische Summe von Spinausrichtungen}.

Jedes Elektron in einem Ferromagneten (Eisen, Nickel, Kobalt)
hat einen Spin --- und damit ein magnetisches Moment.
Die \textbf{Austauschwechselwirkung} (Dokument~174,
Abschnitt über Methode~4) zwingt benachbarte Spins
zur parallelen Ausrichtung: alle up, oder alle down.
Das ist energetisch günstiger als zufällige Verteilung.

Diese kollektive Ausrichtung bildet \textbf{Weiss-Domänen} ---
makroskopische Bereiche kohärenter Spinordnung.
Im Permanentmagneten sind diese Domänen durch ein
äußeres Feld dauerhaft ausgerichtet und eingefroren.

\begin{keyresult}[Was der Permanentmagnet direkt beweist]
\textbf{1. Spin hat bevorzugte, stabile Ausrichtungen.}\\
Up und Down sind nicht gleichwertige Alternativen ---
sie sind geometrisch ausgezeichnete Zustände.
Das Feld \emph{will} in eine von ihnen.

\textbf{2. Diese Stabilität ist makroskopisch robust.}\\
Ein Kühlschrankmagnet hält seine Ausrichtung
bei Raumtemperatur, in einer Umgebung mit $10^{23}$
thermisch fluktuierenden Nachbaratomen ---
über Jahre, ohne zu dekohärieren.

\textbf{3. Superposition tritt kollektiv nie auf.}\\
Kein Permanentmagnet zeigt je einen Zustand
der gleichzeitig Nord und Süd ist.
Die Natur wählt. Immer.
\end{keyresult}

% ============================================================
\section{Die Standardtheorie und ihre Erklärung}
\label{sec:standard}
% ============================================================

Die Standardtheorie erklärt das makroskopische Fehlen
von Superposition mit \textbf{Dekohärenz}:
Die Wechselwirkung mit der Umgebung zerstört die Superposition
so schnell, dass sie praktisch nicht beobachtbar ist.

Das ist technisch korrekt --- aber es ist eine
\textbf{nachträgliche Erklärung}, keine Vorhersage.

\begin{important}[Die strukturelle Schwäche der Dekohärenz-Erklärung]
Die Standardtheorie postuliert:
\begin{enumerate}
  \item Ein einzelnes Elektron existiert in echter Superposition.
  \item Dekohärenz durch die Umgebung zerstört diese Superposition
    praktisch sofort für makroskopische Systeme.
\end{enumerate}

Damit erklärt sie zwar \emph{warum} wir keine makroskopische
Superposition sehen --- aber sie kehrt die physikalische
Reihenfolge um:

\textbf{Sie beginnt mit Superposition als Grundzustand}
und erklärt dann warum dieser Grundzustand immer sofort
zerstört wird.

\textbf{Die FFGFT beginnt umgekehrt:}
Stabilität ist der Grundzustand.
Superposition ist die begrenzte, fragile Ausnahme ---
mit messbaren Zeitskalen und einer geometrischen Ursache.

Die Frage ist nicht nur mathematisch --- sie ist physikalisch:
Welche Beschreibung ist die natürlichere?
\end{important}

% ============================================================
\section{Superposition ist fragil --- die experimentellen Belege}
\label{sec:fragilitaet}
% ============================================================

Die gemessenen Kohärenzzeiten aus Dokument~174
erzählen eine eindeutige Geschichte:

\begin{center}
\resizebox{\textwidth}{!}{%
\begin{tabular}{lrrrl}
\toprule
\textbf{System} & \textbf{$T_1$} & \textbf{$T_2$} &
\textbf{Temperatur} & \textbf{Bemerkung} \\
\midrule
Exziton (CdSe-QD)    & $\sim 1$\,ns  & $\sim 10$\,ps  & 4\,K  &
  Superposition hält Pikosekunden \\
Elektron-Spin (GaAs) & $\sim 1$\,ms  & $\sim 10$\,ns  & mK    &
  Superposition hält Nanosekunden \\
Elektron-Spin (Si)   & $> 1$\,s      & $\sim 10$\,ms  & mK    &
  Superposition hält Millisekunden \\
Kernspin (P:Si)      & $> 1$\,s      & $> 100$\,ms    & mK    &
  Bestes bekanntes System \\
\textbf{Permanentmagnet} & \textbf{Jahre} & \textbf{---} &
  \textbf{300\,K} & \textbf{Stabilität ohne Kühlung} \\
\bottomrule
\end{tabular}
}
\end{center}

Die Tabelle zeigt eine klare Tendenz:
Je isolierter das System vom störenden Einfluss der Umgebung,
desto länger hält die Superposition ---
aber sie hält \emph{immer} nur begrenzt.

\begin{keypoint}[Superposition ist nie der stabile Grundzustand]
Kein einziges experimentelles System zeigt
Superposition als dauerhaft stabilen Zustand.
Selbst im besten Qubit-System (Kernspin P:Si bei mK)
zerfällt die Superposition nach Sekunden.

Der Permanentmagnet --- ohne Kühlung, ohne Isolation ---
hält seinen stabilen Spinzustand über Jahre.

Das ist kein Zufall. Das ist Physik:
\textbf{Stabilität ist energetisch bevorzugt.
Superposition ist metastabil.}
\end{keypoint}

% ============================================================
\section{Das Einzelelektron: Superposition oder unbekannter Zustand?}
\label{sec:einzelelektron}
% ============================================================

Hier liegt der eigentliche konzeptuelle Kern.

Die Standardtheorie behauptet: Ein einzelnes, isoliertes Elektron
befindet sich in echter Superposition
$\alpha\ket{\uparrow} + \beta\ket{\downarrow}$
bis zur Messung.

Die FFGFT-Position ist präziser:

\begin{foundation}[FFGFT: Superposition vs.\ epistemische Unkenntnis]
Ein einzelnes Elektron hat zu jedem Zeitpunkt eine
\textbf{definite Torsionskonfiguration} ---
entweder Spin-up oder Spin-down.

Was wir als Superposition $\alpha\ket{\uparrow} + \beta\ket{\downarrow}$
beschreiben, ist \textbf{epistemische Unkenntnis} über diesen
definitiven Zustand --- keine ontologische Gleichzeitigkeit.

Der Beweis kommt aus der QND-Messung (Dokument~174,
Faraday-Rotation): Ein Spinzustand kann wiederholt gemessen werden,
ohne ihn zu verändern, und liefert \emph{immer dasselbe Ergebnis}.
Das ist nur möglich wenn der Zustand \emph{vor der Messung}
definit existiert.

\textbf{Wenn das Elektron wirklich in Superposition wäre,}
würde eine QND-Messung bei jeder Wiederholung zufällige Ergebnisse
liefern. Das tut sie nicht.
\end{foundation}

Die scheinbare Superposition entsteht in der Beschreibung ---
wenn wir ein präpariertes System vor der Auslesewechselwirkung
mathematisch beschreiben.
Sie ist eine Zustandsbeschreibung unter Unkenntnis,
kein physikalischer Doppelzustand.

% ============================================================
\section{Die Brücke: vom Einzelspin zum Permanentmagneten}
\label{sec:bruecke}
% ============================================================

Der Übergang vom Quanten- zum klassischen Magnetismus
ist in FFGFT kein Rätsel --- er ist derselbe Mechanismus,
nur auf verschiedenen Skalen.

\begin{center}
\resizebox{\textwidth}{!}{%
\begin{tabular}{p{3cm}p{4cm}p{4cm}p{3cm}}
\toprule
\textbf{Ebene} & \textbf{System} & \textbf{FFGFT-Beschreibung}
& \textbf{Stabilität} \\
\midrule
Einzelspin &
  Elektron im QD &
  Eine stabile Torsionskonfiguration; up oder down &
  $T_1 > 1$\,s (P:Si) \\[6pt]
Spin-Paar &
  Zwei gekoppelte QDs &
  Zwei Torsionen; Austauschwechselwirkung bevorzugt Parallelausrichtung &
  $J(t)$ steuerbar \\[6pt]
Domäne &
  Weiss-Domäne ($10^{15}$ Spins) &
  Kollektive Torsionsordnung; energetisch stabiler als Unordnung &
  Stabil bei 300\,K \\[6pt]
\textbf{Permanentmagnet} &
  Makroskopischer Körper &
  Eingefrorene globale Torsionsausrichtung &
  \textbf{Stabil über Jahre} \\
\bottomrule
\end{tabular}
}
\end{center}

Jede Stufe ist dieselbe Physik --- stabile Torsionskonfigurationen,
kollektiv verstärkt durch Austauschwechselwirkung.
Kein Schritt in dieser Kette erfordert echte Superposition.
Jeder Schritt erfordert \textbf{geometrische Präzision und
energetische Stabilität}.

\begin{foundation}[FFGFT: Permanentmagnet als makroskopischer Resonator]
In FFGFT ist der Permanentmagnet kein Sonderfall ---
er ist der natürliche makroskopische Grenzfall
des Resonatorprinzips aus Dokument~173.

Ein einzelner Quantenpunkt auf $\xi$-Stufe $N \approx 8$:
zwei stabile Torsionskonfigurationen, metastabile Zwischenzustände.

Ein Permanentmagnet:
$\sim 10^{22}$ Resonatoren die durch Austauschwechselwirkung
in dieselbe stabile Torsionskonfiguration gezwungen werden.
Das kollektive Feld verstärkt die Stabilität --- statt sie zu
zerstören (Dekohärenz) wirkt die Kopplung hier als
Stabilitätsverstärker.

Die Grenze zwischen Quantenmechanik und klassischem Magnetismus
ist in FFGFT keine kategoriale Grenze ---
es ist ein quantitativer Übergang entlang der $\xi$-Hierarchie.
\end{foundation}

% ============================================================
\section{Was die Standardtheorie schuldig bleibt}
\label{sec:schuldig}
% ============================================================

Die Standarddarstellung von Spin in Lehrbüchern und
populärwissenschaftlichen Texten hat ein strukturelles Problem:
Sie beginnt mit dem Abstrakten (Hilbertraum, Spinor, Kollaps)
und erklärt dann den Alltag (Permanentmagnet) als Grenzfall.

Dabei übersieht sie eine einfache Beobachtung:

\begin{significance}[Die einfachste Beschreibung des Spins]
Der Permanentmagnet existiert.
Er zeigt, dass Spin eine \textbf{räumliche Ausrichtung} ist ---
diskret, stabil, real.

Das ist die anschaulichste Darstellung des Spins:
nicht ein abstrakter Zustandsvektor im Hilbertraum,
sondern eine physikalische Ausrichtung im Raum
die makroskopisch messbar, stabil und alltäglich ist.

Die Quantisierung (nur up oder down, nichts dazwischen)
folgt aus der Geometrie des Resonators --- wie in Dokument~173
hergeleitet. Sie ist kein Postulat, sondern geometrische Konsequenz.

Die scheinbare Superposition ist epistemische Beschreibung
unter Unkenntnis --- kein physikalischer Doppelzustand.
Der Permanentmagnet zeigt täglich, welcher Zustand fundamental ist:
der stabile, ausgerichtete --- nicht der superponierte.
\end{significance}

% ============================================================
\section{Konsequenzen für die Quanteninformation}
\label{sec:qi-konsequenzen}
% ============================================================

Diese Sichtweise hat direkte Konsequenzen für das Verständnis
von Quantencomputern.

Die Standardsicht: Ein Qubit ist mächtig weil es in Superposition ist.
Dekohärenz ist der Feind der Superposition und muss bekämpft werden.

Die FFGFT-Sicht: Ein Qubit ist mächtig weil es eine
\textbf{präzise steuerbare Torsionskonfiguration} ist.
Dekohärenz ist keine Zerstörung eines Wunders ---
sie ist das System das in seinen energetisch stabilen
Grundzustand zurückkehrt. Das ist normal, nicht pathologisch.

\begin{keypoint}[Dekohärenz als Rückkehr zum Grundzustand]
In FFGFT ist Dekohärenz kein Versagen des Systems ---
es ist der Beweis dass der stabile Spinzustand real und
energetisch bevorzugt ist.

Der Permanentmagnet dekohäriert nicht.
Das Qubit dekohäriert in Millisekunden.
Der Unterschied ist nicht die Physik ---
der Unterschied ist die \textbf{Isolation von der Kopplung}
die die stabile Ausrichtung erzwingt.

Ein Quantencomputer kämpft nicht gegen die Natur des Spins ---
er kämpft gegen die Tendenz des Spins zu seiner natürlichen
Stabilität zurückzukehren, bevor die Rechnung fertig ist.
Das ist ein Ingenieursproblem, kein Grundlagenproblem.
\end{keypoint}

% ============================================================
\section{Zusammenfassung}
\label{sec:zusammenfassung-178}
% ============================================================

\begin{conclusionBox}[Fazit: Der Permanentmagnet als Zeuge]
Der Permanentmagnet ist das einfachste und überzeugendste
Argument für die FFGFT-Sicht auf den Spin:

\textbf{1. Spin hat stabile, bevorzugte Ausrichtungen.}\\
Up und Down sind geometrisch ausgezeichnete Zustände ---
nicht willkürliche Messergebnisse.

\textbf{2. Kollektive Superposition existiert nicht.}\\
Milliarden von Elektronen wählen immer eine Ausrichtung.
Niemals beides gleichzeitig. Niemals.

\textbf{3. Superposition ist fragil --- Stabilität ist robust.}\\
Kohärenzzeiten reichen von Pikosekunden (Exziton) bis
Sekunden (Kernspin) --- immer begrenzt, immer endlich.
Der Permanentmagnet hält seinen Zustand Jahre.

\textbf{4. Die Dekohärenz-Erklärung kehrt die Physik um.}\\
Die Standardtheorie beginnt mit Superposition als Grundzustand
und erklärt dann warum er immer sofort verschwindet.
FFGFT beginnt mit Stabilität als Grundzustand ---
Superposition ist die endliche, messbar fragile Ausnahme.

\textbf{5. QND-Messung bestätigt Determinismus.}\\
Wiederholte Faraday-Messungen liefern immer dasselbe Ergebnis ---
der Spinzustand war \emph{vor} der Messung definit.

In FFGFT-Sprache:
Spin-up und Spin-down sind zwei stabile Torsionskonfigurationen.
Der Permanentmagnet ist deren makroskopische kollektive Realisierung.
Superposition ist eine metastabile Zwischentorsion ---
real auf kurzen Zeitskalen, aber nie der fundamentale Grundzustand.

\textbf{Die einfachste Beschreibung des Spins beginnt nicht im
Hilbertraum --- sie beginnt am Kühlschrank.}
\end{conclusionBox}

% ============================================================

\begin{thebibliography}{99}

% --- Permanentmagnetismus und Ferromagnetismus ---
\bibitem{Heisenberg1928}
Heisenberg, W.
\textit{Zur Theorie des Ferromagnetismus.}
Z.\ Phys.\ \textbf{49}, 619--636 (1928).

\bibitem{Weiss1907}
Weiss, P.
\textit{L'hypothèse du champ moléculaire et la propriété ferromagnétique.}
J.\ Phys.\ Théor.\ Appl.\ \textbf{6}, 661--690 (1907).

% --- Austauschwechselwirkung und Spin-Gitter ---
\bibitem{Loss1998}
Loss, D.\ \& DiVincenzo, D.\,P.
\textit{Quantum computation with quantum dots.}
Phys.\ Rev.\ A \textbf{57}, 120--126 (1998).

\bibitem{Petta2005}
Petta, J.\,R.\ et al.
\textit{Coherent manipulation of coupled electron spins in semiconductor
quantum dots.}
Science \textbf{309}, 2180--2184 (2005).

% --- Dekohärenz und Kohärenzzeiten ---
\bibitem{Veldhorst2014}
Veldhorst, M.\ et al.
\textit{An addressable quantum dot qubit with fault-tolerant control-fidelity.}
Nature Nanotechnol.\ \textbf{9}, 981--985 (2014).

\bibitem{Muhonen2014}
Muhonen, J.\,T.\ et al.
\textit{Storing quantum information for 30 seconds in a nanoelectronic device.}
Nature Nanotechnol.\ \textbf{9}, 986--991 (2014).

% --- QND-Messung ---
\bibitem{Atature2007}
Atatüre, M.\ et al.
\textit{Observation of Faraday rotation from a single confined spin.}
Nature Phys.\ \textbf{3}, 101--105 (2007).

\bibitem{Delbecq2019}
Delbecq, M.\,R.\ et al.
\textit{Quantum non-demolition measurement of an electron spin qubit.}
Nature Nanotechnol.\ \textbf{14}, 446--450 (2019).

% --- Dekohärenz-Theorie ---
\bibitem{Zurek2003}
Zurek, W.\,H.
\textit{Decoherence, einselection, and the quantum origins of the classical.}
Rev.\ Mod.\ Phys.\ \textbf{75}, 715--775 (2003).

\bibitem{Joos2003}
Joos, E.\ et al.
\textit{Decoherence and the Appearance of a Classical World in Quantum Theory.}
2nd ed.\ Springer (2003).

% --- Quantenmessung und Determinismus ---
\bibitem{Grangier1998}
Grangier, P., Levenson, J.\,A.\ \& Poizat, J.\,P.
\textit{Quantum non-demolition measurements in optics.}
Nature \textbf{396}, 537--542 (1998).

\end{thebibliography}
